% ============================================================================
% Hub Cultural — Master's Thesis (Mémoire de Master)
% Outline / skeleton — LaTeX, standard academic report structure
% Language: English. Generated as a starting skeleton; content notes are
% marked with \todo{...} and should be replaced with real prose chapter by
% chapter. Numbers/facts already known from the project are inserted as
% placeholders to anchor the writing — verify each before final submission.
% ============================================================================

\documentclass[12pt, a4paper]{report}

% ── Packages ────────────────────────────────────────────────────────────
\usepackage[utf8]{inputenc}
\usepackage[T1]{fontenc}
\usepackage[english]{babel}
\usepackage[margin=2.5cm]{geometry}
\usepackage{graphicx}
\usepackage{hyperref}
\usepackage{natbib}          % swap for biblatex if your university requires it
\usepackage{listings}
\usepackage{xcolor}
\usepackage{amsmath}
\usepackage{booktabs}
\usepackage{caption}
\usepackage{setspace}
\onehalfspacing

% Lightweight TODO marker (deliberately not using the todonotes package —
% avoids a tikz/pgf version dependency that may not match every LaTeX
% install). Remove or search-replace \todo{...} once each section is
% actually written.
\newcommand{\todo}[1]{\par\noindent\textcolor{red}{\textbf{[TODO:}} #1
  \textcolor{red}{\textbf{]}}\par\vspace{0.3em}}

\lstset{basicstyle=\ttfamily\small, breaklines=true, frame=single}

\title{
    % Working title — scope is deliberately broad: cultural relevance, not
    % national origin, is the inclusion criterion (see \S1.1 and \S6.2).
    \textbf{Hub Cultural: Mapping Cultural Event Networks in France \\
    through Instagram Graph Mining and Network Analysis} \\
    \vspace{0.5cm}
    \large From a Diaspora Network to a General Cultural Discovery Platform
}
\author{Diego Merchán}
\date{2026}

% ============================================================================
\begin{document}

\maketitle

\begin{abstract}
\todo{150--250 words. Draft last, once Results/Discussion are stable.
Should state: (1) the problem — culturally-relevant activity in the Paris
region (with priority given to photography, cinema, literature, theatre,
and local community events) is not systematically documented or
discoverable, and this project originated from, but is no longer limited
to, the specific case of the Latin American diaspora; (2) the method — a
pipeline that scrapes Instagram, builds a Neo4j graph, applies network
analysis (PageRank, betweenness, Leiden community detection) and an
LLM/NLP event-extraction pipeline, combined with human curation as the
trust signal, plus a research-operations discipline (decision log, run
log, dry-run-first design) treated as methodology in its own right; (3)
headline findings — network size and structure, the bridge role of
commercial/institutional accounts, event landscape, and a transparent
account of data-quality limits; (4) the practical artifact — a public,
bilingual event-discovery website; (5) contribution — methodological
(human curation vs.\ automated scoring; a cultural-relevance criterion
replacing a nationality-based one) and empirical (first mapped network of
this kind for this cultural ecosystem in France).}
\end{abstract}

\tableofcontents

% ============================================================================
\chapter{Introduction}
% ============================================================================

\section{Context and Motivation}

Hub Cultural began as an attempt to make visible something that is well
known anecdotally but poorly documented empirically: the cultural life
that Latin American, and in particular Colombian, communities sustain in
the Paris region through small associations, independent venues, cultural
attachés, and a dense but scattered presence on Instagram. The author's
own Colombian background supplied both the initial motivation and the
first seed account for data collection — the Colombian consulate in
Paris (\texttt{@consuladocolparis}) — from which the network was grown
outward through automatic relationship discovery.

As the project progressed, the boundary that had organized it from the
start — \emph{is this account or event part of the Latin American
diaspora?} — proved both practically brittle and conceptually narrower
than what the data, and the underlying research interest, actually
supported. Many of the best-connected cultural accounts discovered
through the network turned out to be French institutions, mixed-nationality
collectives, or venues with no diaspora affiliation at all, but which were
unmistakably part of the same cultural ecosystem the project set out to
map — programming photography exhibitions, independent cinema, literary
events, and theatre alongside explicitly Colombian or Latin American
programming. The project's scope therefore evolved in two steps: first
from Colombia-specific to Latin-American-diaspora-wide, and subsequently
from a nationality-based criterion to a purely cultural one. The operative
question became not \emph{who is behind this account} but \emph{is this
culturally relevant}, with explicit priority — inherited from the
project's original four thematic agendas (\texttt{docs/objetivos.md}) —
given to photography, cinema, literature, and theatre, alongside a fifth,
less formally bounded category: events organized by and for local
communities, regardless of national or cultural origin.

This evolution is treated in this thesis not as scope creep to be
apologized for, but as a methodological finding in its own right
(\S\ref{sec:scope-evolution}): a nationality-based inclusion criterion
turned out to be both harder to operationalize consistently and less
faithful to how cultural networks actually cohere in practice, where
diaspora and host-society cultural production overlap and reinforce each
other rather than remaining cleanly separable.

\section{Problem Statement and Research Question}

The cultural activity produced by, and adjacent to, migrant and diaspora
communities in a host country is rarely centralized: it lives across
dozens of Instagram accounts, is announced inconsistently, and is legible
mainly to whoever already follows the right accounts. This is true of
Latin American cultural life in the Paris region specifically, but — as
this project's own data eventually made clear — it is not a problem
unique to any one national community; it is a structural feature of how
small-scale cultural production circulates on social media generally.

This thesis asks: \emph{How is culturally-relevant activity — with
priority given to photography, cinema, literature, theatre, and local
community cultural life — organized as a network of accounts and events
in the Paris region, and can that structure be systematically mapped,
analyzed, and made discoverable using publicly available Instagram
data?}

This breaks into three sub-questions:
\begin{enumerate}
    \item \textbf{Structural.} Who are the connective actors in this
    network — which accounts function as hubs or bridges between
    otherwise separate cultural circles, and does this differ between
    commercial/institutional accounts and individual or grassroots ones?
    \item \textbf{Thematic.} How is cultural activity distributed across
    the project's priority domains (photography, cinema, literature,
    theatre, community events), and how well does an automated pipeline
    classify it into these domains compared to human judgment?
    \item \textbf{Practical.} Can the resulting graph be turned into a
    public-facing tool that helps the community it describes actually
    discover and attend these events — and what does building that tool
    reveal about the reliability of the underlying data?
\end{enumerate}

\section{Objectives}

The project's objectives, as originally framed (\texttt{docs/objetivos.md})
and updated here to reflect the current, cultural-relevance-based scope:

\begin{enumerate}
    \item Build and maintain four thematic monitoring agendas —
    \emph{literature}, \emph{photography}, \emph{cinema}, and
    \emph{theatre} — tracking the organizations that play an active role
    in each, complemented by a fifth, less formally bounded priority:
    cultural events organized by and for local communities.
    \item Consolidate a pipeline that extracts events from the accounts
    under these agendas, primarily from Instagram and, prospectively,
    from websites and other social networks.
    \item Map the network of actors within each thematic agenda and
    categorize them — by type of organization, geographic zone, cultural
    identity where relevant, and structural role in the network.
    \item Evaluate the participation opportunities these networks open
    up, using a deliberately human, non-technical success criterion: on
    this dimension, the project succeeds when it leads to actually
    meeting someone through it.
\end{enumerate}

\section{Scope and Contributions}

\textbf{Scope.} Geographically the project centers on Île-de-France,
though a meaningful share of curated accounts and events originate from,
or are organized by, institutions physically located outside France — a
scope boundary discussed explicitly in \S\ref{sec:limitations} rather
than silently discarded. Instagram is the sole data source, a deliberate
boundary rather than an oversight: it is the platform where this cultural
activity is, in practice, announced. Only public accounts are in scope.
Inclusion is governed by cultural relevance rather than national origin,
with the four thematic pillars and the community-events category acting
as explicit priorities rather than hard filters.

\textbf{Contributions.} This thesis makes three kinds of contribution:
\begin{itemize}
    \item \textit{Methodological} — a working pipeline combining network
    analysis, multilingual NLP event extraction, and human curation as
    the operative trust signal, arrived at through an evidenced pivot
    away from fully automated account scoring
    (\S\ref{sec:methodology-curation}); and a research-operations
    practice — a running decision log, a run log, and a dry-run-first
    discipline applied to every script — treated here as part of the
    methodology rather than incidental bookkeeping
    (\S\ref{sec:research-ops}).
    \item \textit{Empirical} — a systematic mapping of this cultural
    network's structure: node and edge counts, community structure, and
    the role of bridging actors, alongside a transparent account of
    where the underlying data is, and is not, reliable (geocoding
    confidence, date-extraction accuracy, categorization accuracy).
    \item \textit{Applied} — a public-facing, bilingual event-discovery
    website that turns the graph into something a member of the
    community it describes can actually use, not only an analytical
    artifact for the researcher.
\end{itemize}

\section{Thesis Structure}

Chapter~2 situates the project within diaspora and migration studies,
social network analysis, computational social science, and NLP-based
event extraction. Chapter~3 details the methodology, including the two
methodological pivots this thesis treats as contributions in their own
right: human curation as the trust signal, and a research-operations
discipline of decision logging, run logging, and dry-run-first
iteration. Chapter~4 describes the resulting system, from data
collection through the public discovery website. Chapter~5 reports
results: network structure, thematic and geographic distribution, and a
transparent accounting of data quality. Chapter~6 discusses what these
results mean, including a dedicated reflection on the project's
evolution from a nationality-based to a cultural-relevance-based scope.
Chapter~7 concludes, states limitations, and outlines future work.


% ============================================================================
\chapter{Literature Review / Theoretical Framework}
% ============================================================================

\section{Diaspora Studies and Cultural Networks}
\todo{Situate within diaspora/migration studies: how migrant
communities organize culturally, the role of associations and informal
networks, prior work on Latin American diasporas in Europe specifically
if available. Framing note: this remains one relevant thread — the
project's motivating case and a meaningful share of its curated content
— but since the operative inclusion criterion is now cultural relevance
rather than national origin (\S\ref{sec:scope-evolution}), pair this with
literature on cultural-sector/creative-industry network mapping more
generally, so the theoretical framing matches the project's actual scope
and not only its origin story.}

\section{Social Network Analysis in Migration Contexts}
\todo{Prior applications of SNA to migrant/diaspora networks;
concepts to define here for later use: centrality (degree, betweenness,
PageRank), community detection, structural holes / bridging, E-I index.}

\section{Computational Social Science and Social Media as Data}
\todo{Justify Instagram as a data source; discuss known biases
(who is/isn't represented on Instagram, algorithmic visibility, platform
ToS and ethical use of public data).}

\section{Graph Databases and Network Science Methods}
\todo{Neo4j / property graph model; why a graph database over a
relational one for this problem; brief technical grounding for PageRank,
betweenness centrality, Leiden algorithm, k-core decomposition,
participation coefficient — cite original papers.}

\section{NLP for Event Extraction and Classification}
\todo{Ground the 3-layer detection design: sentence embeddings
for candidate filtering, NLI zero-shot classification, LLM-based
extraction/classification/gating. Cite relevant NLP literature (semantic
similarity, zero-shot NLI, LLM information extraction) and note the
multilingual requirement (ES/FR/EN) as a specific design constraint.}

\section{Positioning}
\todo{One paragraph: how this thesis differs from prior diaspora
network studies (data-driven, automated + human-curated hybrid, produces
a live public tool, not just an analysis) and from prior event-extraction
NLP work (domain-specific gating logic, human-curation-as-trust-signal
pivot).}


% ============================================================================
\chapter{Methodology}
\label{ch:methodology}
% ============================================================================

\section{Research Design Overview}
\todo{High-level pipeline diagram description (7 phases):
extraction $\to$ ingestion $\to$ network analysis $\to$ manual
categorization $\to$ event enrichment $\to$ geo enrichment $\to$
cleanup. Figure: reproduce/adapt the pipeline diagram from
\texttt{docs/data\_flow.md}.}

\section{Research Operations: Decision Logging, Run Tracking, and Iterative Safety}
\label{sec:research-ops}

Independently of the technical pipeline itself, this project maintained
three interlocking practices for managing the research process, treated
here as a methodological contribution rather than incidental project
management. Together they make the pipeline's evolution auditable after
the fact — which is not free in a setting where the object of study is a
continuously growing graph, reprocessed repeatedly by successive versions
of the extraction logic.

\textbf{A real-time decision log.} Every non-trivial design choice was
recorded, at the moment it was made, in a sequentially numbered decision
log (\texttt{docs/decisions.md}, entries DD-001 through DD-044 and
counting) rather than reconstructed retrospectively for this thesis. Each
entry states the alternatives considered, the evidence that motivated the
final choice, and, where applicable, validation results. Three entries
illustrate the range of what this captured: DD-041 documents a
geo-conflict guardrail added to the event-deduplication resolver after a
production dry-run surfaced a specific failure mode — two structurally
similar events merged despite being in different cities — validated
against a 14-case test suite before being applied to 68 pending merges;
DD-042 documents the addition of a per-event, LLM-generated tag field
designed specifically to avoid a parsing bug present in an earlier,
account-level free-text field; DD-043 and DD-044 document the design of
the Neo4j-to-static-JSON export and the resulting public website,
including an explicit accounting of which ranking components could be
precomputed server-side versus which needed to remain client-side.
Maintaining this log in real time, rather than writing it up after the
fact, is itself a methodological choice: it captures the reasoning
available \emph{at the time} of each decision, including alternatives
that were considered and rejected — something retrospective documentation
systematically loses.

\textbf{A run log.} Separately, every substantive execution of the
pipeline against live data — not every script invocation, but every run
whose output was retained or acted upon — was logged with its date,
parameters, and a summary of results (\texttt{docs/runs\_log.md}). This
served two purposes beyond record-keeping: it made it possible to
distinguish a genuine change in the underlying data (new accounts or
events discovered) from a change in pipeline logic when comparing two
exports, and it created a paper trail for decisions that depended on a
specific run's output rather than on memory — for instance, the decision
not to add a venue-specificity guardrail to the event resolver was made
with explicit reference to a logged run's results.

\textbf{Dry-run-first, diagnose-first design.} Every script that writes
to the live Neo4j graph was built with a \texttt{--dry-run} or
\texttt{--diagnose} mode as a first-class feature, not an afterthought:
\texttt{cleanup\_legacy\_accounts.py} runs its deletion logic inside a
single transaction and issues a \texttt{ROLLBACK} under
\texttt{--dry-run}, so a preview's reported counts are exact rather than
estimated; \texttt{load\_manual\_account\_categorization.py} and the
event-extraction and event-resolution scripts all follow the same
pattern. Combined with the idempotent design of the extraction and
ingestion scripts — re-running them is always safe, since already-
processed accounts and posts are skipped or merged rather than
duplicated — this made it possible to iterate on pipeline logic against
a live, shared, non-trivially-sized research database without the usual
risk of an irreversible mistake, a constraint that shaped several of the
decisions recorded in the log above.

Together, these three practices form a coherent research-operations
discipline: decisions are recorded when made, runs are recorded when
executed, and changes to the live data are always previewable before
being committed. For a thesis built on a continuously evolving dataset
rather than a fixed one, this discipline is what makes it possible to
state, for any given result in Chapter~\ref{ch:results}, exactly which
version of the pipeline and which run produced it.

\section{Data Collection}
\todo{Apify actors used (instagram-profile-scraper,
instagram-post-scraper); seed selection strategy — originally
\texttt{@consuladocolparis} as single seed, BFS-style expansion via
\texttt{RELATED\_TO} suggestions, later curated seed files
(\texttt{config/seeds\_*.json}); cost control mechanism
(\texttt{.apify\_cost\_log.json}); incremental/idempotent design
(\texttt{extract\_profiles.py} only scrapes accounts not yet seen).}

\section{Data Model}
\todo{Full Neo4j schema: node labels (\texttt{:Account},
\texttt{:Post}, \texttt{:Hashtag}, \texttt{:Location}, \texttt{:Track},
\texttt{:Comment}, \texttt{:IgtvVideo}, \texttt{:Event}, \texttt{:City},
\texttt{:Country}, \texttt{:Arrondissement}) and relationship types
(\texttt{PUBLISHED}, \texttt{HAS\_HASHTAG}, \texttt{MENTIONS},
\texttt{TAGS\_USER}, \texttt{COAUTHORED\_BY}, \texttt{MENTIONS\_EVENT},
\texttt{ORGANIZED}, \texttt{PARTICIPATED\_IN}, \texttt{LOCATED\_IN},
etc.). Include an entity-relationship diagram (Appendix~\ref{app:schema}).}

\section{Network Analysis Methods}
\todo{Multiplex network design: social layer (MENTIONS/
TAGS\_USER/COAUTHORED\_BY, author$\to$post$\to$target projection)
separate from the algorithmic layer (\texttt{RELATED\_TO}, \texttt{Algo}
suffix). Metrics computed: exact PageRank, exact betweenness, Leiden
community detection at three resolutions ($\gamma = 0.5, 1.0, 1.5$,
seed 42), weakly connected components, k-core, participation
coefficient, E-I index by actor type. Note the pivot from Neo4j GDS
(blocked — standard AuraDB lacks the plugin, corporate firewall also
blocks port 7687) to a local igraph/leidenalg implementation
(\texttt{3\_analyze\_network.py}) as a methodological adaptation worth
documenting, not hiding.}

\section{Event Detection Pipeline}
\todo{Describe the 3-layer architecture as it exists today
(post the 2026-08 rewrite that removed the old embeddings-based
multi-label typification layer after a 100-post human-vs-pipeline
evaluation — cite \texttt{data\_processed/eval\_100\_report.md}):
\begin{enumerate}
    \item Layer 1 — cosine similarity against $\sim$100 reference
    phrases (\texttt{paraphrase-multilingual-MiniLM-L12-v2}) as a
    candidate filter.
    \item Layer 2 — binary NLI classification
    (\texttt{MoritzLaurer/multilingual-MiniLMv2-L6-mnli-xnli}),
    hypothesis picked in the caption's detected language.
    \item Layer 3 — LLM (Ollama/Groq/Cerebras) performs event typing
    (16-label taxonomy), price-range extraction, and critically
    \emph{gates} event creation based on
    \texttt{is\_public\_invitation}/\texttt{is\_upcoming}.
\end{enumerate}
Report the false-positive reduction this gating produced, per the
100-post eval.}

\section{Human Curation as the Trust Signal}
\label{sec:methodology-curation}
\todo{This is a genuine methodological pivot worth a dedicated
section: the original automated account-scoring classifier
(\texttt{old/1\_harvest\_account\_classifier.py}) selected accounts
where well under 10\% turned out to be culturally valuable on manual
review. \texttt{manualDataCuratedAt} (set only by
\texttt{load\_manual\_account\_categorization.py}) replaced the
scoring-based selection as the operative "worth keeping" signal.
Discuss this as a finding in itself: automated relevance scoring
under-performed simple human judgment for this domain, and why (cite
whatever intuition emerged — commercial/political accounts scoring
artificially high, genuinely small cultural accounts scoring low on
network centrality despite being sociologically important).}

\section{Geospatial Enrichment}
\todo{Nominatim geocoding, rate-limited (1.1s/req); hierarchy
\texttt{Location}$\to$\texttt{Arrondissement}$\to$\texttt{City}$\to$
\texttt{Country}; three-tier confidence strategy (exact / city-combined
/ city-hint-only) and — \textbf{report honestly} — the empirical
finding that low-confidence geocodes need explicit exclusion rather
than fallback display, since a fallback to a default city hint can
silently produce plausible-looking but wrong coordinates (found during
dashboard validation, \S\ref{sec:limitations}).}

\section{Ranking / Relevance Scoring}
\label{sec:ranking}
\todo{Document the discovery-website ranking formula as a
methodological artifact in its own right:
$\text{Relevance} = 100 \cdot (0.30Q + 0.22A + 0.18B + 0.20T + 0.10C)
\cdot \text{penalty}$, where $Q$ is extraction quality/confidence,
$A$ is network centrality (percentile-blended PageRank/betweenness/
followers/participation), $B$ is social resonance (hotness/post
count), $T$ is temporal proximity, $C$ is session-based personalization
with no login. Note the explicit removal of a nationality/
"Colombian-ness" bonus once the project's scope broadened — a
methodological choice worth justifying (avoiding a nationality bias
baked into a supposedly neutral relevance score).}

\section{Ethical Considerations}
\todo{Public-data-only scope; Instagram ToS; no scraping of
private accounts; anonymization/aggregation choices for anyone
appearing in captions/comments who is not an organizational account;
data retention and the decision to eventually delete non-curated
accounts (\texttt{cleanup\_legacy\_accounts.py}).}

\section{Limitations of Method}
\label{sec:limitations}
\todo{Be concrete and honest — this section should draw on
real bugs found during the dashboard validation, not hypothetical
caveats: (1) geocoding fallback producing false-precision coordinates
when a location could not be reliably resolved; (2) LLM date
extraction misreading year/date from ambiguous captions (e.g.\ a
caption stating only a month, defaulting to the wrong year); (3)
Instagram-only sourcing missing organizations that primarily use other
channels; (4) accounts outside the intended geographic scope
(e.g.\ France-branded institutions physically located in Colombia)
occasionally entering the graph via automatic relationship discovery
before being caught by curation.}


% ============================================================================
\chapter{System Implementation: The Hub Cultural Platform}
% ============================================================================

\section{Pipeline Architecture Overview}
\todo{Reproduce the phase diagram; emphasize idempotency and
incrementality as design principles running through every script.}

\section{Extraction and Ingestion}
\todo{\texttt{1\_harvest\_ig\_profiles.py} /
\texttt{1\_harvest\_ig\_posts.py} / \texttt{extract\_profiles.py} $\to$
\texttt{2\_build\_graph.py}. Provenance tracking
(\texttt{firstSeenAt}/\texttt{lastUpdatedAt}).}

\section{Network Analysis Engine}
\todo{\texttt{3\_analyze\_network.py} internals; why it runs
entirely offline from exported CSVs (portability, reproducibility
benefit worth noting for a thesis — the analysis can be re-run without
live database access).}

\section{The Discovery Website}
\todo{Describe the client-facing artifact built from
\texttt{5\_export\_dashboard\_data.py} $\to$ static site (\texttt{site/}):
dynamic category menu (a deliberate choice — no fixed editorial
category list, so the menu reflects whatever the upstream human
curation actually prioritizes), bilingual ES/FR interface, no-login
session personalization via \texttt{localStorage}, similarity search
via embeddings. Frame this as the thesis's applied contribution:
turning the graph into something a member of the community could
actually use, not just an analytical artifact for the researcher.}

\section{Decision Log as Methodological Transparency}
\todo{Note the project's practice of a running, dated decision
log (\texttt{docs/decisions.md}, DD-001 through DD-044+) as a form of
methodological transparency — cite it as supporting material and
summarize its most consequential entries (the curation pivot, the
GDS$\to$igraph pivot, the event-pipeline rewrite, the geocoding
confidence gating fix).}


% ============================================================================
\chapter{Results}
\label{ch:results}
% ============================================================================

\section{Network Structure}
\todo{Report current numbers precisely — re-pull from the
latest \texttt{data\_processed/*.csv} before writing, do not reuse V1
numbers unchanged (V1 reference point: $\sim$7{,}200 nodes seeded from
a single account, 61 communities via Leiden, high modularity). State
current account/event/location counts as of the most recent run
(cf.\ latest export: 170 events, 200 accounts, at time of dashboard
validation — confirm this is post-cleanup and reflects the broadened
non-Colombia-only scope before citing).}

\section{Bridge Actors and Community Structure}
\todo{The V1 sociological finding — commercial/institutional
accounts show high E-I index (structural bridges) while individual
accounts cluster locally — needs re-validation against the current
graph (broadened scope, curation changes since V1). Report whether it
still holds.}

\section{Thematic and Event Landscape}
\todo{Distribution of events by category (musical, visual,
institutional, festival, academic, etc.); how this compares to the
four originally-targeted thematic agendas (literature, photography,
cinema, theatre); note the empirical finding that
\texttt{gastronomía}/institutional content over-represents relative to
the curated thematic priorities, and how the dashboard's category
menu and curation approach responded to that.}

\section{Geographic Distribution}
\todo{Coverage: total events vs.\ events with a resolvable
\texttt{Location} node vs.\ events with real lat/lon (report the
actual, current, post-fix percentage — do not reuse the pre-fix
85.5\% figure without re-checking after the geocoding-confidence
exclusion fix described in \S\ref{sec:limitations}).}

\section{Data Quality Findings}
\todo{Treat this as a genuine result, not just a limitations
footnote: quantify how often the geocoding fallback produced
false-precision coordinates before the fix, how often LLM date
extraction required manual correction, what fraction of institutional-
category events were miscategorized on manual sample review (cf.\ the
20-event institutional sample: $\sim$35\% either miscategorized or
pipeline-gate leaks). This kind of transparent error analysis is
strong thesis material.}


% ============================================================================
\chapter{Discussion}
\label{ch:discussion}
% ============================================================================

\section{Interpreting Network Structure}
\todo{What does the bridge-actor pattern (if it holds) say
sociologically about how this diaspora organizes culturally —
mediated through a small number of institutional hubs rather than a
dense peer network?}

\section{From a Diaspora Network to a Cultural-Relevance Network}
\label{sec:scope-evolution}

The project's scope changed twice. It began as a network of Colombian
diaspora accounts, seeded from a single institutional account, the
Colombian consulate in Paris. As the network grew through automatic
relationship discovery (\texttt{RELATED\_TO}), it broadened to Latin
American diaspora accounts more generally. Both of these were still
nationality-based criteria: an account or event qualified because of
\emph{who} was behind it.

The second, more consequential change replaced that criterion with a
purely cultural one: an account or event qualifies if it is culturally
relevant, with explicit priority given to photography, cinema,
literature, theatre, and locally-organized community events, regardless
of the nationality or cultural origin of the organizers behind it. This
is a stricter methodological change than it might first appear, because
it also required removing a nationality bonus that had previously been
present in the discovery website's relevance-ranking formula
(\S\ref{sec:ranking}) — an explicit multiplicative bonus applied to
events tagged with a Colombian cultural identity. Retaining that bonus
after the scope change would have produced a ranking signal that was
formally neutral (it ranks \emph{cultural} events) while materially
still favoring one nationality over every other one present in the same
graph; removing it was necessary for the ranking formula to be honest
about what it actually measures.

\todo{Once network composition is re-pulled for Chapter~\ref{ch:results},
quantify here: what share of curated accounts/events are Colombian,
other Latin American, French, or other, before vs.\ after the scope
change; and revisit whether the four thematic pillars (literature,
photography, cinema, theatre) still describe the bulk of curated content
or whether the fifth, community-events category has grown large enough
to warrant its own treatment in the results and in the discovery
website's category menu.}

What this change gained, beyond ranking-formula honesty, is a criterion
that is easier to defend and to operationalize consistently:
\emph{is this cultural, and does it fall within one of these priority
domains} is a judgment that can be applied uniformly to any account the
discovery process surfaces. \emph{Is this part of the Latin American
diaspora}, by contrast, required an ad hoc call for every borderline
case — a French venue that regularly hosts Colombian artists, a Latin
American cultural institute whose visitors are mostly French — that the
pipeline had no principled way to resolve consistently. What it cost is
some of the project's original sociological specificity: the initial
research question was, in part, about how \emph{one} diaspora organizes
culturally in a host country, and that question is necessarily diluted
once the object of study is cultural activity in general. This thesis
handles that trade-off by keeping the diaspora-specific reading as one
lens applied to a subset of the network (wherever \texttt{culturalIdentity}
data supports it), rather than as the network's defining boundary.

\section{On Data Quality and the Limits of Automation}
\todo{Tie together \S\ref{sec:methodology-curation} and
\S\ref{sec:limitations}: this project's most consistent lesson is that
fully automated signals (account-relevance scoring, geocoding
fallbacks, LLM date extraction) degrade silently and look confident
even when wrong, while human curation — despite being unscalable — was
the more reliable trust signal throughout. Discuss what this implies
for similar future projects.}

\section{Practical Implications: The Website as Intervention}
\todo{Beyond analysis, the discovery website is itself a
sociological intervention — it could actually change how the diaspora
finds/attends events. Discuss this as a design-research contribution,
distinct from the analytical contribution.}

\section{Methodological Reflections}
\todo{What would be done differently starting over — e.g.\
building the geocoding-confidence gate and human curation trust signal
in from day one rather than arriving at them through iteration.}


% ============================================================================
\chapter{Conclusion}
\label{ch:conclusion}
% ============================================================================

\section{Summary of Contributions}
\todo{Restate the three contributions from \S1.4 concretely
against what was actually delivered.}

\section{Answering the Research Question}
\todo{Direct answer, one paragraph.}

\section{Limitations}
\todo{Consolidate — do not just repeat \S\ref{sec:limitations}
verbatim; frame as boundaries on how far the findings generalize.}

\section{Future Work}
\todo{From the project's own V2 planning notes: a new NLP
account classifier with explicit geography/cultural-relevance
dimensions and anti-embeddings to penalize commercial/political
content; expanded seed set from Latin American consulates in Paris;
revised harvest cycle prioritizing cheap profile scraping before
expensive post scraping; remaining dashboard work (clickable
arrondissement/commune map, dedicated Explorer/Map/Organizer-profile
pages, per-event bilingual content generated at extraction time,
generic category image set).}

\section{Closing Reflection}
\todo{Personal note tying back to \S1.3's fourth objective —
whether/how the project succeeded by its own most human metric
(meeting someone through it), separate from the technical metrics.}


% ============================================================================
\appendix
% ============================================================================

\chapter{Graph Schema Reference}
\label{app:schema}
\todo{Full node/relationship table, adapted from
\texttt{docs/data\_flow.md} and \texttt{CLAUDE.md}.}

\chapter{Selected Decision Log Entries}
\todo{Reproduce/summarize the most thesis-relevant entries
from \texttt{docs/decisions.md} (curation pivot, GDS$\to$igraph pivot,
event-pipeline rewrite \& 100-post eval, geocoding confidence gate).}

\chapter{Discovery Website — Screenshots}
\todo{Home screen, filter bar, event detail panel — capture
once the site is stable post-fixes (see current punch list: upcoming-
filter bug, geocoding exclusion, account-scope cleanup).}

\chapter{Repository Structure}
\todo{Reproduce the script naming convention and phase
breakdown from \texttt{CLAUDE.md}'s Commands section as a reference
table.}

% ============================================================================
% \bibliographystyle{apalike}   % or whatever citation style your program requires
% \bibliography{references}
% ============================================================================

\end{document}
